{\small\textbf{Back to March 2020 (Zheng, 1:30) [23:00–24:30]}}\par\smallskip
Let me answer the three questions of the assignment.
\begin{itemize}\setlength\itemsep{0pt}
\item \textbf{The environment:} a recovered economy hit by war, inflation and rate hikes; credit cold; housing hot, then cooling; strong, profitable banks.
\item \textbf{The criteria and risks:} a two-percent normal-times level, calibrated on past losses, insuring against shocks no one can forecast. Housing went to other tools, and structural buffers were cut in return.
\item \textbf{The implications:} cheap, largely a swap, no sign of contraction, complementary to monetary policy, reciprocated abroad, and 6.7 billion euros now releasable. The price is reliance on discretion.
\end{itemize}
\par\smallskip
If there is one idea to take away, it is this: the Dutch decision was less about how \emph{much} capital banks hold, and more about \emph{which kind}. Capital that the authority can release when the next shock comes, whatever its source.
\par\smallskip
\emph{(pause)} In March 2020, the Netherlands had nothing to release. \textbf{Next time, it will.}
\par\smallskip
The best time to fill a buffer is when nobody thinks you need it, even in a storm. Thank you; we look forward to your questions.
